% ---------------------------------------------------------------------------
% Submission to Physical Review D (APS).  REVTeX 4.2.
% Compile: pdflatex a0_half_dark_energy_rate.tex  (twice, for references)
% ---------------------------------------------------------------------------
\documentclass[aps,prd,reprint,amsmath,amssymb,nofootinbib,superscriptaddress]{revtex4-2}

\usepackage{graphicx}
\usepackage{bm}

\begin{document}

\title{The MOND acceleration scale is one half of the dark-energy gravitational rate}

\author{Carl P. Zimmerman}
\affiliation{Briar Creek Tech, Charlotte, North Carolina, USA}

\date{\today}

\begin{abstract}
The acceleration scale $a_0$ that organizes galactic rotation curves is known to lie near $cH_0$.
We record an exact closed form for it in terms of the dark-energy density alone,
$a_0 = \tfrac{1}{2}c\sqrt{G\rho_\Lambda}$, equivalently $c^2\sqrt{\Lambda/32\pi}$.
With the Planck 2018 $\rho_\Lambda$ this gives $9.43\times10^{-11}\,\mathrm{m\,s^{-2}}$, which is
$0.70\%$ from the value $9.36\times10^{-11}\,\mathrm{m\,s^{-2}}$ obtained by fitting 175 SPARC
rotation curves at a stellar mass-to-light ratio $\Upsilon_{[3.6]}=0.70$ (residual scatter
$0.108$~dex). The expression is parameter-free apart from the rational coefficient $\tfrac12$.
We state the limits of the claim explicitly. The interpolation used is identical to that of
Milgrom [Phys.\ Lett.\ A \textbf{253}, 273 (1999)], which also \emph{derives}
$a_0 = 2cH_\Lambda$, a factor $11.58$ larger; and the present coefficient is not observationally
distinguishable from Milgrom's later empirical $cH_\Lambda/2\pi$, the two differing by $7.9\%$
($0.018$~dex in $g_{\rm obs}$, $16\%$ of the fit scatter). The contribution is therefore one of
\emph{form}---an exact expression in $\rho_\Lambda$ with a rational coefficient---and not of improved
fit. The coefficient is fitted, not derived. We give one falsifiable consequence: if
$a_0\propto\sqrt{\rho_\Lambda}$ then $a_0(z)$ is constant for $w=-1$ and rises to a maximum before
declining for the DESI-preferred $w_0>-1$, $w_a<0$, in contrast to the monotonic rise predicted by
$a_0 \propto cH(z)$.
\end{abstract}

\maketitle

\section{The relation}

Let $\rho_\Lambda \equiv \Lambda c^2/8\pi G$ be the dark-energy mass density and let
$\sqrt{G\rho_\Lambda}$ denote the associated gravitational rate, the inverse of the dynamical time of a
region of that density. The relation recorded here is
\begin{equation}
a_0 \;=\; \tfrac{1}{2}\,c\sqrt{G\rho_\Lambda}
      \;=\; c^{2}\sqrt{\frac{\Lambda}{32\pi}}
      \;=\; \frac{cH_\Lambda}{\sqrt{32\pi/3}} ,
\label{eq:main}
\end{equation}
where $H_\Lambda = c\sqrt{\Lambda/3} = \sqrt{8\pi G\rho_\Lambda/3}$. The three forms are algebraically
identical. The reduction is immediate and worth displaying, because it shows that the factor
$\sqrt{32\pi/3}$ carries no structural content:
\begin{equation}
\frac{cH_\Lambda}{\sqrt{32\pi/3}}
= c\sqrt{\frac{8\pi G\rho_\Lambda}{3}}\sqrt{\frac{3}{32\pi}}
= c\sqrt{\frac{G\rho_\Lambda}{4}}
= \tfrac12 c\sqrt{G\rho_\Lambda}.
\label{eq:reduce}
\end{equation}
Every factor of $\pi$ cancels, as does the $3$. Accordingly $\sqrt{32\pi/3}=5.78881$ should be read as
the numerical conversion between $H_\Lambda$ and $\sqrt{G\rho_\Lambda}$ bookkeeping, and not as a
geometric quantity to be explained. The entire content of Eq.~(\ref{eq:main}) is the single rational
coefficient $\tfrac12$.

\section{Relation to prior work}

We place this before the numerical evaluation, since it bounds what can be claimed.

Milgrom introduced $a_0$ and noted its proximity to $cH_0$~\cite{Milgrom1983}. The interpolation used
throughout the present work,
\begin{equation}
g_{\rm obs}^2 = g_{\rm bar}^2 + a_0\,g_{\rm bar},
\qquad
\nu(y) = \sqrt{1+1/y},
\label{eq:closure}
\end{equation}
with $y = g_{\rm bar}/a_0$, is \emph{identical} to Eq.~(9) of Ref.~\cite{Milgrom1999}; it is not a
variant of it. That same reference derives a value for the scale from a de Sitter--Unruh vacuum
argument and obtains $a_0 = 2cH_\Lambda$, which exceeds Eq.~(\ref{eq:main}) by the factor
$2\sqrt{32\pi/3} = 11.58$. The de Sitter--Unruh route therefore does not leave the coefficient
undetermined: it predicts one, and that prediction misses the measured scale by an order of magnitude.
The present contribution is a renormalization of that coefficient to match data, not a new derivation.
Comparable $2cH$-type coefficients are obtained in Refs.~\cite{Pikhitsa2010,Klinkhamer2011}.

The sharpest comparison is with the empirical $a_0 = cH_\Lambda/2\pi$ of Ref.~\cite{Milgrom2020}. Since
$2\pi = 6.28319$ against $\sqrt{32\pi/3} = 5.78881$, the two coefficients differ by $7.9\%$, i.e.\
$0.036$~dex in $a_0$. In the low-acceleration regime $g_{\rm obs}\simeq\sqrt{g_{\rm bar}a_0}$ this
enters $g_{\rm obs}$ at half that, $0.018$~dex, against a fit scatter of $0.108$~dex
(Sec.~\ref{sec:num}). \emph{The two coefficients are observationally indistinguishable at present
precision.} We therefore make no claim of improved fit. What distinguishes Eq.~(\ref{eq:main}) is that
it is a closed form in $\rho_\Lambda$ with a rational coefficient, where $cH_\Lambda/2\pi$ is an
empirical divisor.

Finally, we note for completeness that modified-inertia realizations of this phenomenology, their
generic nonlocality, and the orbit-shape dependence of their effective interpolating function are
established in Refs.~\cite{Milgrom1994,Milgrom2022}; no result from that line is claimed here.

\section{Numerical evaluation}
\label{sec:num}

Using Planck 2018 TT,TE,EE+lowE+lensing+BAO values $H_0 = 67.66\ \mathrm{km\,s^{-1}\,Mpc^{-1}}$ and
$\Omega_\Lambda = 0.6889$~\cite{Planck2018} gives
$\rho_\Lambda = 5.924\times10^{-27}\,\mathrm{kg\,m^{-3}}$, $H_\Lambda = 1.820\times10^{-18}\,\mathrm{s^{-1}}$,
and
\begin{equation}
a_0 = \tfrac12 c\sqrt{G\rho_\Lambda} = 9.43\times10^{-11}\ \mathrm{m\,s^{-2}} .
\end{equation}

Fitting Eq.~(\ref{eq:closure}) to the 175 galaxies of the SPARC sample~\cite{SPARC} with $a_0$ held at
this value and a single global stellar mass-to-light ratio as the only free parameter returns
$\Upsilon_{\rm disk} = 0.70$ (with $\Upsilon_{\rm bulge} = 1.4\,\Upsilon_{\rm disk} = 0.98$) at a
residual radial-acceleration-relation scatter of $0.108$~dex. For reference, the same fit with the
conventional $a_0 = 1.2\times10^{-10}\,\mathrm{m\,s^{-2}}$ and $\Upsilon_{\rm disk} = 0.5$ gives
$0.122$~dex. We stress that $a_0$ and $\Upsilon$ are degenerate in this fit, so the comparison should
not be read as a measurement of $a_0$; the relevant statement is that
$\Upsilon_{\rm disk} = 0.70$ lies inside the range $0.5$--$0.8$ expected for Spitzer $3.6\,\mu$m
stellar populations, so Eq.~(\ref{eq:main}) does not require an implausible mass-to-light ratio.

One alternative normalization must be reported. Using the total density together with $cH_0$ in place of
the pure-$\Lambda$ combination yields $a_0 = 1.13\times10^{-10}\,\mathrm{m\,s^{-2}}$, $21\%$ larger.
The present fit does not discriminate between these; Sec.~\ref{sec:test} gives an observable that does.

\section{What is not claimed}

The coefficient $\tfrac12$ is fitted. We attempted to fix it from ghost-freedom, unitarity, and a
holographic degree-of-freedom count and did not succeed: the constructions that reproduce
Eq.~(\ref{eq:closure}) leave the overall normalization free, and the one geometric candidate that
appeared to fix it, $(3/8\pi)^{1/4} = 0.5878$, is the single-degree-of-freedom limit of a count with no
independent microscopic justification. Equation~(\ref{eq:main}) should therefore be read as a
prediction of the \emph{scaling} $a_0\propto\sqrt{\rho_\Lambda}$ with an unexplained coefficient, and
not as an explanation of the value of $a_0$. No claim is made here regarding particle physics or
unification.

\section{A falsifiable consequence}
\label{sec:test}

The two normalizations of Sec.~\ref{sec:num} differ in redshift dependence. If
$a_0\propto\sqrt{\rho_\Lambda}$ with a dark-energy equation of state
$w(z) = w_0 + w_a z/(1+z)$, then
\begin{equation}
\frac{a_0(z)}{a_0(0)} = (1+z)^{\frac32(1+w_0+w_a)}\exp\!\left[-\frac{3w_a z}{2(1+z)}\right].
\label{eq:evol}
\end{equation}
For $w_0=-1$, $w_a=0$ this is exactly constant. For the DESI-preferred region $w_0>-1$, $w_a<0$ it
rises to a maximum and then declines; it is not a monotonic rise. By contrast $a_0\propto cH(z)$ rises
monotonically. Measurements of the acceleration scale at $z\gtrsim1$ therefore discriminate between the
normalizations, and a monotonic rise would exclude Eq.~(\ref{eq:main}).

We note that at least one recent determination reports a rising $a_0(z)$, which is in tension with
Eq.~(\ref{eq:evol}) on the pure-$\Lambda$ branch. That measurement is itself degenerate with
$\Lambda$CDM assumptions in the mass modeling, and we record it as a live tension rather than as
support.

\section{Known tensions}

We list these so that they are not overlooked.

\emph{Clusters.} Equation~(\ref{eq:closure}) applied at cluster scales underpredicts the required
acceleration. On uncorrected eRASS1 masses the median discrepancy factor at $R_{500}$ is $2.33$, and the
cluster-scale acceleration inferred from lensing, $2.0\times10^{-9}\,\mathrm{m\,s^{-2}}$, exceeds
Eq.~(\ref{eq:main}) by a factor $21.6$. A weak-lensing mass recalibration reduces the former to roughly
$1.6$--$1.8$ but does not remove it. Because Eq.~(\ref{eq:main}) gives a \emph{lower} $a_0$ than the
conventional value, this discrepancy is $13\%$ worse here than for the conventional coefficient. The
problem is shared with relativistic MOND theories generally and is unresolved.

\emph{Solar system.} Applied exactly, Eq.~(\ref{eq:closure}) leaves a constant residual acceleration
$a_0/2$ that exceeds Earth--Mars ranging bounds by three orders of magnitude. This is a property of the
specific interpolation, not of Eq.~(\ref{eq:main}): an interpolation with a faster Newtonian approach,
$\mu(x) = x/\sqrt{1+x^2}$, leaves a residual $a_0^2/2g$ instead and reduces the anomaly by five orders
of magnitude, at a cost of $0.003$~dex in the SPARC fit. The word ``exact'' should not be attached to
Eq.~(\ref{eq:closure}).

\emph{Degeneracy.} As Sec.~\ref{sec:num} notes, the rotation-curve data constrain the combination of
$a_0$ and $\Upsilon$, not $a_0$ alone. Nothing here should be read as determining $a_0$ to better than
the $7.9\%$ separating Eq.~(\ref{eq:main}) from $cH_\Lambda/2\pi$.

\section{Methods and provenance}

The relation Eq.~(\ref{eq:main}) was arrived at independently by the author through three routes, which
are recorded here for transparency about how the result was obtained.

\begin{enumerate}
\item An agentic automated-research harness~\cite{AutoresearchTool}, adapted by the author to
      search candidate dimensional combinations of $c$, $G$, $\Lambda$ and $H_0$ against the measured
      scale. (The cited harness is general-purpose infrastructure for running automated research agents;
      the search task applied here is the author's own.)
\item Symbolic regression over those dimensional combinations, in the spirit of the AI Feynman
      method of Udrescu and Tegmark~\cite{AIFeynman}, which recovered
      $a_0 \propto c\sqrt{G\rho_\Lambda}$ as the lowest-complexity expression consistent with the
      target value.
\item Direct fitting of Eq.~(\ref{eq:closure}) to the SPARC rotation-curve
      compilation~\cite{SPARC}, which fixes the coefficient at $\tfrac12$ to the precision quoted in
      Sec.~\ref{sec:num}.
\end{enumerate}

Routes 1 and 2 identified the functional form; route 3 fixed the coefficient. We emphasize that none of
the three constitutes a derivation, and that a coefficient obtained by fitting remains a fitted
coefficient however it was found (Sec.~IV). All numerical claims in this paper are reproduced by
publicly available scripts that exit with a nonzero status if any internal consistency check fails
\cite{Repo}.

\section*{Disclosure of AI assistance}

Portions of the analysis, the numerical verification, and the drafting of this paper were carried out
with the assistance of a large language model (Anthropic Claude). The author directed the work,
specified and reviewed every load-bearing calculation, and takes full responsibility for the contents,
including any errors. No artificial-intelligence system satisfies the criteria for authorship and none
is listed as an author. Several intermediate claims produced during the work were subsequently found to
be incorrect and were withdrawn prior to submission; the claims presented here are those that survived
independent re-derivation.

\begin{acknowledgments}
The author thanks the SPARC team for making the rotation-curve compilation publicly available, and
acknowledges the AI Feynman project for the symbolic-regression approach adopted in Sec.~VII.
\end{acknowledgments}

\begin{thebibliography}{99}

\bibitem{Milgrom1983} M.~Milgrom, Astrophys.\ J.\ \textbf{270}, 365 (1983).

\bibitem{Milgrom1999} M.~Milgrom, Phys.\ Lett.\ A \textbf{253}, 273 (1999).

\bibitem{Pikhitsa2010} P.~V.~Pikhitsa, arXiv:1010.0318.

\bibitem{Klinkhamer2011} F.~R.~Klinkhamer and M.~Kopp, arXiv:1104.2022.

\bibitem{Milgrom2020} M.~Milgrom, Phys.\ Rev.\ D \textbf{102}, 084010 (2020).

\bibitem{Milgrom1994} M.~Milgrom, Ann.\ Phys.\ (N.Y.) \textbf{229}, 384 (1994).

\bibitem{Milgrom2022} M.~Milgrom, Phys.\ Rev.\ D \textbf{106}, 064060 (2022).

\bibitem{Planck2018} Planck Collaboration, Astron.\ Astrophys.\ \textbf{641}, A6 (2020).

\bibitem{SPARC} F.~Lelli, S.~S.~McGaugh, and J.~M.~Schombert, Astron.\ J.\ \textbf{152}, 157 (2016).

\bibitem{AIFeynman} S.-M.~Udrescu and M.~Tegmark, Sci.\ Adv.\ \textbf{6}, eaay2631 (2020).

\bibitem{AutoresearchTool} A.~Karpathy, \emph{autoresearch}, software repository (MIT license),
  \url{https://github.com/karpathy/autoresearch}.

\bibitem{Repo} C.~P.~Zimmerman, analysis repository,
  \url{https://github.com/carlzimmerman/zimmerman-formula}.

\end{thebibliography}

\end{document}
